\documentclass[12pt,a4paper]{article}

\usepackage[utf8]{inputenc}
\usepackage[T1]{fontenc}
\usepackage{lmodern}
\usepackage{amsmath,amssymb,amsthm}
\usepackage{geometry}
\usepackage{setspace}
\usepackage{booktabs}
\usepackage{tabularx}
\usepackage{natbib}
\usepackage{hyperref}
\usepackage{enumitem}

\geometry{margin=1in}
\onehalfspacing

\newtheorem{proposition}{Proposition}
\newtheorem{corollary}{Corollary}
\newtheorem{assumption}{Assumption}
\newtheorem{remark}{Remark}

\title{\textbf{Debt Denomination, Fiscal Backing, and Inflation:\\ Evidence from Brazilian Treasury Debt}}
\author{Research Proposal}
\date{\today}

\begin{document}

\maketitle

\begin{abstract}
This paper studies whether the denomination of government liabilities shapes the transmission of fiscal shocks to inflation and asset prices. Fiscal valuation theories imply that inflation can reduce the real burden of nominal fixed-rate debt, while inflation-indexed liabilities cannot be revalued through the price level. Floating-rate debt occupies an intermediate position because monetary tightening offsets part of the fiscal relief produced by inflation. These contractual differences generate an inflation-ranking restriction across Brazilian Treasury instruments: nominal fixed-rate debt shocks should generate stronger inflation responses than floating-rate shocks, which should in turn dominate inflation-indexed shocks.

The central empirical challenge is that debt issuance shocks are not necessarily fiscal valuation shocks. A debt expansion need not be inflationary if agents simultaneously revise upward expected future primary surpluses. We therefore construct \emph{survey-purged debt issuance shocks}: issuance innovations orthogonal to revisions in fiscal backing measured from the Prisma Fiscal survey. We show that debt denomination generates an inflation-ranking restriction that is robust to a nontrivial region of residual fiscal-backing heterogeneity. The primary empirical test compares nominal and indexed debt, whose valuation elasticities are pinned down by contract. The sharper test exploits the full nominal--floating--indexed ranking.

Using monthly Brazilian data on Treasury issuance, IPCA inflation, SELIC rates, indexed bond yields, sovereign spreads, and fiscal expectations, we estimate local projections for inflation and valuation variables. The paper contributes a new identification strategy for fiscal theories of inflation by using debt denomination to obtain observable restrictions that do not require direct measurement of infinite-horizon fiscal backing.
\end{abstract}

\noindent \textbf{JEL codes:} E31, E62, H63, E52.\\
\textbf{Keywords:} Fiscal Theory of the Price Level, debt denomination, inflation-indexed debt, fiscal backing, local projections, Brazil.

\section{Introduction}

Understanding the relationship between public debt and inflation remains a central question in macroeconomics. The Fiscal Theory of the Price Level (FTPL) provides a valuation-based view of price determination: the real value of government liabilities must equal the expected present value of future primary surpluses. When fiscal backing is insufficient, the price level may adjust to reduce the real value of outstanding nominal obligations.

Despite its theoretical clarity, empirical evaluation of fiscal valuation mechanisms is difficult. The relevant object, the expected present value of future primary surpluses, is forward looking and only imperfectly observable. Inflation following a fiscal expansion may be interpreted as evidence of insufficient fiscal backing, while the absence of inflation may be rationalized by expectations of future taxes or spending cuts. This observational equivalence has limited direct empirical tests of fiscal theories of inflation.

This paper proposes a different empirical strategy. Rather than attempting to observe infinite-horizon fiscal backing directly, we exploit a contractual feature of government liabilities: debt denomination. Inflation reduces the real value of nominal fixed-rate debt, but it does not reduce the real value of inflation-indexed debt. Floating-rate debt lies between these extremes because inflation tends to induce higher short-term nominal rates, increasing future debt service and partially offsetting the fiscal relief from inflation. Hence, debt denomination generates a ranking of inflation incentives that can be tested in the data.

Brazil provides an unusually rich setting for this test. Since the early 2000s, the National Treasury has actively issued three broad classes of domestic securities: nominal fixed-rate debt (LTN and NTN-F), floating-rate debt indexed to the SELIC rate (LFT), and inflation-indexed debt linked to the IPCA (NTN-B). This instrument menu allows us to compare shocks to liabilities with different valuation elasticities within a common fiscal, monetary, and institutional environment.

The paper makes three contributions. First, it develops a debt-denomination restriction for fiscal valuation. The primary implication is that, for otherwise comparable survey-purged issuance shocks, nominal fixed-rate debt should generate stronger inflation responses than indexed debt. A sharper implication ranks all three instruments: nominal fixed-rate debt should generate the strongest inflation response, floating-rate debt an intermediate response, and indexed debt the weakest response.

Second, the paper introduces survey-purged debt issuance shocks. We first extract issuance innovations for each debt type from forecasting equations based on lagged macro-financial information. We then purge these innovations of contemporaneous revisions in expected fiscal backing using Prisma Fiscal survey expectations. This procedure does not claim to identify perfectly unbacked fiscal shocks. Instead, it constructs shocks that are not rationalized by observable fiscal-backing news.

Third, the paper clarifies the identification content of debt denomination under imperfectly observed backing. We show that the ranking of inflation responses depends on debt instruments' valuation elasticities and survives a nontrivial region of heterogeneity in residual unobserved backing. This makes the test more robust than a direct attempt to measure the level of expected future primary surpluses.

The remainder of the proposal is organized as follows. Section \ref{sec:theory} develops the valuation logic and the ranking restriction. Section \ref{sec:identification} introduces survey-purged shocks and the residual-backing robustness region. Section \ref{sec:data} describes the data. Section \ref{sec:empirical} presents the empirical strategy. Section \ref{sec:robustness} discusses robustness exercises and external validation. Section \ref{sec:conclusion} concludes.

\section{Theory: Debt Denomination and Valuation Elasticities}
\label{sec:theory}

\subsection{Government valuation with multiple debt instruments}

Let \(\mathcal{F}_t\) denote the expected present value of future primary surpluses:
\begin{equation}
\mathcal{F}_t
=
E_t\left[\sum_{h=0}^{\infty} M_{t,t+h} S_{t+h}\right],
\label{eq:fiscal_backing}
\end{equation}
where \(M_{t,t+h}\) is the stochastic discount factor and \(S_{t+h}\) is the real primary surplus.

The government issues three types of domestic liabilities:
\begin{itemize}[leftmargin=1.2cm]
    \item \(B_t^N\): nominal fixed-rate debt (LTN and NTN-F);
    \item \(B_t^F\): floating-rate debt indexed to the SELIC rate (LFT);
    \item \(B_t^R\): inflation-indexed debt linked to the IPCA (NTN-B).
\end{itemize}

Let \(V_t^j\) denote the real market value of debt type \(j\in\{N,F,R\}\). The fiscal valuation equation can be written as
\begin{equation}
V_t^N+V_t^F+V_t^R=\mathcal{F}_t.
\label{eq:valuation_equation}
\end{equation}

For a one-period fixed-rate nominal bond, \(V_t^N=B_t^N/P_t\). For indexed debt, the real payoff is protected against inflation, so its real value does not fall mechanically with \(P_t\). Floating-rate debt is nominal in principal, but its coupons reset with the short-term policy rate. Under an inflation-targeting regime, inflation increases the policy rate and therefore raises future debt service, offsetting part of the valuation relief from inflation.

\subsection{Valuation elasticity}

Define the valuation elasticity of instrument \(j\) with respect to the price level as
\begin{equation}
\kappa^j
\equiv
-\frac{\partial \log V_t^j}{\partial \log P_t},
\label{eq:kappa_def}
\end{equation}
where the derivative is interpreted as the net valuation effect, including induced future interest-rate costs when relevant.

For nominal fixed-rate debt, inflation reduces the real burden one-for-one:
\begin{equation}
\kappa^N=1.
\end{equation}

For inflation-indexed debt, inflation does not reduce the real contractual burden:
\begin{equation}
\kappa^R=0.
\end{equation}

For floating-rate debt, the valuation elasticity is intermediate:
\begin{equation}
0<\kappa^F<1.
\end{equation}
The lower bound reflects the fact that floating-rate debt is still a nominal liability in real terms at the instant of revaluation. The upper bound reflects the monetary-policy feedback: higher inflation raises the SELIC rate and increases future debt-service costs. The exact value of \(\kappa^F\) depends on the monetary regime, especially the strength of the policy-rate response to inflation. The paper does not require estimating \(\kappa^F\); it only uses the ranking
\begin{equation}
\kappa^N>\kappa^F>\kappa^R.
\label{eq:kappa_ranking}
\end{equation}

\subsection{Fiscal valuation versus resource-constraint mechanisms}

Let \(\varepsilon_t^j\) denote a debt issuance shock of type \(j\), after removing observable fiscal-backing news as described below. Under a fiscal-valuation mechanism, the inflation response to a shock is increasing in the valuation elasticity of the affected debt instrument. In reduced form, the valuation component of the inflation response can be written as
\begin{equation}
IRF_{\pi}^{j}(h)
=
\kappa^j \, \Omega_t^j \, \Theta_{\pi}(h),
\label{eq:irf_theory_general}
\end{equation}
where \(\Omega_t^j\) is the fraction of the shock not absorbed by fiscal backing or non-price valuation margins, and \(\Theta_{\pi}(h)>0\) is a common propagation term.

By contrast, resource-constraint theories emphasize aggregate demand relative to productive capacity. Under such mechanisms, debt denomination should be secondary once the size of the fiscal impulse and economic slack are controlled for. Conditional on slack, inflation responses should be similar across debt types:
\begin{equation}
|IRF_{\pi}^{N}(h)|
\approx
|IRF_{\pi}^{F}(h)|
\approx
|IRF_{\pi}^{R}(h)|.
\label{eq:resource_restriction}
\end{equation}

The paper therefore distinguishes a valuation mechanism from a resource-constraint mechanism using cross-instrument rankings. The primary valuation prediction is
\begin{equation}
|IRF_{\pi}^{N}(h)|>|IRF_{\pi}^{R}(h)|,
\label{eq:primary_ranking}
\end{equation}
while the sharper three-instrument prediction is
\begin{equation}
|IRF_{\pi}^{N}(h)|>|IRF_{\pi}^{F}(h)|>|IRF_{\pi}^{R}(h)|.
\label{eq:three_way_ranking}
\end{equation}

Equation \eqref{eq:primary_ranking} compares two contractually extreme instruments: fixed-rate nominal debt has maximal valuation elasticity, while indexed debt has zero direct valuation elasticity. Equation \eqref{eq:three_way_ranking} is a sharper overidentifying implication because it uses the intermediate position of floating-rate debt, whose elasticity depends on the monetary-policy regime.

\section{Identification: Survey-Purged Debt Shocks}
\label{sec:identification}

\subsection{Why issuance surprises are not enough}

A simple debt issuance forecast error is not sufficient to test fiscal valuation. Under FTPL, an unexpected debt expansion need not be inflationary if agents simultaneously expect future primary surpluses to rise. The relevant shock is not issuance per se, but issuance not backed by revisions in expected fiscal surpluses.

Let \(I_t^j\) denote net issuance of debt type \(j\in\{N,F,R\}\), measured as a share of GDP. We first estimate issuance forecasting equations:
\begin{equation}
I_t^j
=
\alpha_j
+
\Lambda_j' Z_{t-1}
+u_t^j,
\qquad
j\in\{N,F,R\},
\label{eq:issuance_forecast}
\end{equation}
where \(Z_{t-1}\) includes lagged macroeconomic variables, fiscal indicators, yield-curve factors, and global risk controls. The residual \(u_t^j\) is an issuance innovation for instrument \(j\).

We then use Prisma Fiscal survey revisions to remove observable fiscal-backing news. Let \(\Delta \mathcal{F}_t^{P}\) denote the Prisma-based revision in expected fiscal backing. The baseline construction is
\begin{equation}
u_t^j
=
\gamma_j \Delta \mathcal{F}_t^{P}
+
\varepsilon_t^j,
\qquad
j\in\{N,F,R\}.
\label{eq:survey_purge}
\end{equation}
The residual \(\varepsilon_t^j\) is a \emph{survey-purged debt issuance shock}. It is an issuance innovation not contemporaneously rationalized by observable revisions in fiscal expectations.

Equation \eqref{eq:survey_purge} is estimated separately for each instrument. There is no theoretical reason to expect \(\gamma_N=\gamma_F=\gamma_R\), and equality of these coefficients is not part of the identification argument. The coefficients \(\gamma_j\) describe how issuance innovations comove with observed Prisma revisions; they do not measure the unobserved residual backing that may remain after the purge.

\subsection{Constructing fiscal-backing revisions}

Let \(E_t^P[s_{t+q}]\) denote the Prisma Fiscal consensus forecast of the primary surplus at horizon \(q\), expressed as a share of GDP. We construct a finite-horizon proxy for fiscal backing:
\begin{equation}
\mathcal{F}_t^{P}
=
\sum_{q=0}^{Q}
\frac{E_t^P[s_{t+q}]}{(1+\bar r)^q},
\label{eq:prisma_pv}
\end{equation}
where \(\bar r\) is a benchmark real discount rate and \(Q\) is determined by the available survey horizon. The revision is
\begin{equation}
\Delta \mathcal{F}_t^{P}
=
\mathcal{F}_t^{P}-\mathcal{F}_{t-1}^{P}.
\label{eq:prisma_revision}
\end{equation}

This measure is not the infinite-horizon object in equation \eqref{eq:fiscal_backing}. Rather, it is a high-frequency proxy for fiscal news. In the Brazilian context, near-term primary-balance revisions are plausibly informative about broader fiscal credibility because fiscal regimes have often been evaluated through compliance with short- and medium-term surplus targets.

\subsection{Residual backing and the invariance theorem}

Survey purging removes observable fiscal-backing news, but it need not remove all backing news. We therefore decompose the remaining response of fiscal backing to a survey-purged shock into an unobserved residual component.

Let \(\rho^j\in[0,1]\) denote the residual backing share for a survey-purged shock in instrument \(j\). A value \(\rho^j=0\) corresponds to a fully unbacked shock after survey purging. A value \(\rho^j=1\) corresponds to a fully backed shock, despite the absence of measured Prisma revisions. The fraction \(1-\rho^j\) is the residual fiscal imbalance that must be absorbed through inflation or other valuation margins.

The valuation component of the inflation response can then be written as
\begin{equation}
IRF_{\pi}^{j}(h)
=
\kappa^j(1-\rho^j)\Theta_{\pi}(h),
\qquad
j\in\{N,F,R\},
\label{eq:irf_rho}
\end{equation}
where \(\Theta_{\pi}(h)>0\) is common across debt types after normalizing shocks to the same size.

\begin{proposition}[Debt-denomination invariance]
Suppose \(\kappa^N>\kappa^F>\kappa^R=0\) and \(\Theta_{\pi}(h)>0\). If residual backing is common across debt instruments, \(\rho^N=\rho^F=\rho^R=\rho<1\), then
\begin{equation}
|IRF_{\pi}^{N}(h)|>|IRF_{\pi}^{F}(h)|>|IRF_{\pi}^{R}(h)|
\end{equation}
for every horizon \(h\) at which \(\Theta_{\pi}(h)>0\). The ranking identifies the relative ordering of valuation elasticities \(\kappa^N,\kappa^F,\kappa^R\), not the level of residual fiscal backing \(\rho\).
\end{proposition}

\begin{proof}
Under common residual backing, equation \eqref{eq:irf_rho} becomes
\[
IRF_{\pi}^{j}(h)=\kappa^j(1-\rho)\Theta_{\pi}(h).
\]
Because \((1-\rho)\Theta_{\pi}(h)>0\), the ordering of inflation responses is the ordering of \(\kappa^j\). Thus \(\kappa^N>\kappa^F>\kappa^R\) implies the stated ranking. The level of \(\rho\) scales all responses equally and therefore cannot be recovered from the ranking. If \(\rho=1\), all valuation-driven inflation responses collapse to zero and the ranking test has no power.\end{proof}

\subsection{Robustness to heterogeneous residual backing}

Common residual backing is sufficient but not necessary. Define
\begin{equation}
\phi^j \equiv \kappa^j(1-\rho^j).
\end{equation}
The three-way inflation ranking survives if and only if
\begin{equation}
\phi^N>\phi^F>\phi^R.
\end{equation}
Equivalently, the admissible residual-backing region is
\begin{align}
\kappa^N\rho^N-\kappa^F\rho^F &< \kappa^N-\kappa^F,
\label{eq:region_nf}\\
\kappa^F\rho^F-\kappa^R\rho^R &< \kappa^F-\kappa^R.
\label{eq:region_fr}
\end{align}
These inequalities define an open convex region in \((\rho^N,\rho^F,\rho^R)\). The common-backing case \(\rho^N=\rho^F=\rho^R=\rho<1\) is a special case within this region.

This characterization clarifies the interpretation of violations. If nominal shocks retain more residual backing than floating or indexed shocks, the ranking is harder to detect because \((1-\rho^N)\) is smaller. In that case, finding \(|IRF_{\pi}^{N}|>|IRF_{\pi}^{F}|>|IRF_{\pi}^{R}|\) is conservative. If instead nominal shocks retain less residual backing than the other instruments, heterogeneous backing may mimic the valuation ranking. The empirical design therefore treats the nominal--indexed comparison as the primary test, because \(\kappa^N=1\) and \(\kappa^R=0\) are contractually pinned down. The full nominal--floating--indexed ranking is a sharper but more regime-sensitive implication.

\section{Data}
\label{sec:data}

The empirical analysis uses monthly Brazilian data. Table \ref{tab:variables} summarizes the main variables.

\begin{table}[h!]
\centering
\caption{Main variables and sources}
\label{tab:variables}
\begin{tabularx}{\textwidth}{lX}
\toprule
Variable & Description and source \\
\midrule
\(I_t^N\) & Net issuance of nominal fixed-rate securities (LTN and NTN-F), Tesouro Nacional. \\
\(I_t^F\) & Net issuance of floating-rate securities (LFT), Tesouro Nacional. \\
\(I_t^R\) & Net issuance of inflation-indexed securities (NTN-B), Tesouro Nacional. \\
\(\pi_t\) & IPCA inflation, IBGE. \\
\(i_t\) & SELIC policy rate, Banco Central do Brasil. \\
\(Y_t\) & IBC-Br economic activity index, Banco Central do Brasil. \\
\(e_t\) & BRL/USD exchange rate, Banco Central do Brasil. \\
\(CDS_t\) & 5-year sovereign CDS spread, Bloomberg or Refinitiv. \\
\(r_t\) & NTN-B real yield by maturity bucket, ANBIMA. \\
\(\mathcal{F}_t^P\) & Prisma Fiscal survey-based fiscal-backing proxy, Ministry of Finance. \\
\bottomrule
\end{tabularx}
\end{table}

The baseline sample starts when the relevant Treasury instruments and Prisma Fiscal expectations are available at sufficient frequency. All issuance variables are normalized by GDP. Inflation, interest rates, real yields, and CDS spreads are transformed consistently across specifications, with robustness checks using alternative transformations.

\section{Empirical Strategy}
\label{sec:empirical}

\subsection{Baseline local projections}

For each outcome \(y_t\) and horizon \(h\), we estimate
\begin{equation}
y_{t+h}-y_{t-1}
=
\alpha_h
+
\beta_h^N\varepsilon_t^N
+
\beta_h^F\varepsilon_t^F
+
\beta_h^R\varepsilon_t^R
+
\Gamma_h'X_t
+\eta_{t+h},
\label{eq:lp_baseline}
\end{equation}
where \(\varepsilon_t^N\), \(\varepsilon_t^F\), and \(\varepsilon_t^R\) are survey-purged debt issuance shocks, and \(X_t\) includes lags of macroeconomic controls, lags of the dependent variable, lags of the shock variables, monetary-policy controls, fiscal controls, and global risk factors.

The main outcomes are:
\begin{enumerate}[leftmargin=1.2cm]
    \item IPCA inflation and inflation expectations;
    \item sovereign CDS spreads;
    \item NTN-B real yields;
    \item SELIC policy rates;
    \item the nominal exchange rate.
\end{enumerate}

The exchange-rate responses are included as auxiliary evidence. Under a valuation mechanism, indexed and floating-rate shocks may exert greater pressure on risk premia and therefore on currency depreciation, while nominal fixed-rate shocks may operate more directly through domestic prices. Because exchange-rate dynamics also reflect global risk and monetary-policy channels, these responses are interpreted cautiously.

\subsection{Primary and sharper tests}

The primary inflation test is
\begin{equation}
H_{V}^{NR}:\quad |\beta_h^{N,\pi}|>|\beta_h^{R,\pi}|.
\label{eq:test_nr}
\end{equation}
This test compares the two contractually extreme valuation elasticities.

The sharper three-way test is
\begin{equation}
H_{V}^{NFR}:\quad |\beta_h^{N,\pi}|>|\beta_h^{F,\pi}|>|\beta_h^{R,\pi}|.
\label{eq:test_nfr}
\end{equation}
This test uses the intermediate valuation elasticity of floating-rate debt and is therefore more informative but more sensitive to the monetary regime.

For valuation variables, the fiscal-valuation mechanism predicts relatively greater pressure from instruments that inflation cannot easily revalue:
\begin{equation}
|\beta_h^{R,z}|\gtrsim |\beta_h^{F,z}|\gtrsim |\beta_h^{N,z}|,
\qquad
z\in\{CDS,r\}.
\label{eq:valuation_variables}
\end{equation}
Equation \eqref{eq:valuation_variables} is interpreted as a directional implication rather than a strict theorem, especially for CDS spreads, which include currency-risk and global-risk components outside the domestic valuation equation.

\subsection{State dependence and resource constraints}

To distinguish valuation mechanisms from resource-constraint mechanisms, we allow responses to vary with economic slack. Let \(S_{t-1}=1\) denote a high-slack state and \(S_{t-1}=0\) denote a near-capacity state. We estimate
\begin{align}
y_{t+h}-y_{t-1}
=&
S_{t-1}
\left(
\beta_{h,s}^{N}\varepsilon_t^N
+
\beta_{h,s}^{F}\varepsilon_t^F
+
\beta_{h,s}^{R}\varepsilon_t^R
\right)
\nonumber\\
&+
(1-S_{t-1})
\left(
\beta_{h,c}^{N}\varepsilon_t^N
+
\beta_{h,c}^{F}\varepsilon_t^F
+
\beta_{h,c}^{R}\varepsilon_t^R
\right)
+
\Gamma_h'X_t
+\eta_{t+h}.
\label{eq:lp_state}
\end{align}

Under a resource-constraint view, inflation responses should depend primarily on slack and should be similar across debt instruments conditional on slack:
\begin{equation}
|\beta_{h,q}^{N,\pi}|
\approx
|\beta_{h,q}^{F,\pi}|
\approx
|\beta_{h,q}^{R,\pi}|,
\qquad
q\in\{s,c\}.
\end{equation}
Under a valuation view, debt-denomination rankings should remain visible within slack and capacity states, although their magnitudes may vary with macroeconomic conditions.

\section{Robustness, External Validation, and Alternative Identification}
\label{sec:robustness}

\subsection{Narrative validation}

We construct a monthly narrative database of fiscal and debt-management events using Treasury auction reports, annual borrowing plans, Ministry of Finance announcements, congressional budget documents, IMF Article IV reports, and rating-agency actions. Events are classified as:
\begin{enumerate}[leftmargin=1.2cm]
    \item unanticipated issuance expansions;
    \item composition shifts across nominal, floating-rate, and indexed debt;
    \item fiscal-backing news that should be absorbed by the survey-purging procedure.
\end{enumerate}

Let \(N_t^{j,unant}\) denote an unanticipated narrative event associated with debt type \(j\). We validate the survey-purged shocks by estimating
\begin{equation}
\varepsilon_t^j
=
\delta_0^j
+
\delta_1^j N_t^{j,unant}
+
\xi_t^j.
\end{equation}
A positive \(\delta_1^j\) supports the interpretation that the statistical shocks correspond to historically meaningful debt-management innovations. We also check cross-type orthogonality by testing whether narrative events for debt type \(k\) predict shocks for debt type \(j\neq k\).

\subsection{Pension-fund regulatory shocks as IV robustness}

As an external validation exercise, we use regulatory changes affecting closed pension funds' demand for NTN-B securities. Brazilian pension funds are natural holders of inflation-indexed bonds because their liabilities are often indexed to inflation. Regulatory changes that alter asset-liability matching requirements shift demand for NTN-B securities for reasons plausibly unrelated to contemporaneous fiscal backing.

Let \(Z_t^R\) denote predicted NTN-B demand induced by regulatory changes. The IV first stage is
\begin{equation}
\varepsilon_t^R
=
\delta Z_t^R
+
\Phi'X_t
+
\nu_t.
\end{equation}
The corresponding IV local projection replaces \(\varepsilon_t^R\) in equation \eqref{eq:lp_baseline} with its fitted component. This IV strategy is treated as robustness rather than the baseline because its exclusion restriction may be threatened by direct effects of regulatory demand on secondary-market NTN-B yields or rollover risk.

We report weak-instrument diagnostics, Anderson--Rubin confidence sets, and sensitivity to excluding individual regulatory events. We also test whether the regulatory instrument predicts outcomes before announcement dates or directly affects secondary-market yields after controlling for indexed issuance.

\subsection{Additional robustness checks}

The paper will report the following robustness exercises:
\begin{enumerate}[leftmargin=1.2cm]
    \item alternative constructions of \(\mathcal{F}_t^P\), including different discount rates and forecast horizons;
    \item alternative issuance forecasting models for equation \eqref{eq:issuance_forecast};
    \item alternative slack measures, including output gaps, unemployment gaps, capacity utilization, and composite labor-market slack indicators;
    \item exclusion of crisis periods, including the global financial crisis, the Brazilian fiscal crisis, and the COVID-19 shock;
    \item placebo tests in which survey-purged shocks predict future revisions in \(\mathcal{F}_{t+k}^{P}\);
    \item alternative sovereign-risk measures, including EMBI spreads and domestic term premia;
    \item machine-learning forecasting models for issuance as robustness only, not as primary identification.
\end{enumerate}

\section{Expected Contributions}

The proposed paper contributes to three literatures. First, it contributes to empirical work on fiscal theories of inflation by providing an observable ranking restriction based on debt denomination. Second, it contributes to the identification of fiscal shocks by combining issuance innovations with survey-based fiscal expectation revisions. Third, it contributes to the literature on sovereign debt management by showing that the composition of liabilities may affect the transmission of fiscal policy to inflation and risk premia.

The key conceptual contribution is that debt denomination generates an inflation-ranking restriction that is robust to unobserved fiscal backing over a nontrivial region of residual-backing heterogeneity. This shifts the empirical focus from estimating the level of fiscal backing to testing a relative prediction across instruments with different valuation elasticities.

\section{Conclusion}
\label{sec:conclusion}

This paper proposes a new empirical approach to fiscal inflation based on the denomination of government debt. The main insight is that nominal fixed-rate, floating-rate, and inflation-indexed liabilities differ in their valuation elasticity with respect to the price level. Fiscal valuation therefore predicts a ranking of inflation responses across debt types.

Using Brazilian Treasury data and Prisma Fiscal expectations, we construct survey-purged issuance shocks that remove observable fiscal-backing news. We then test whether inflation and valuation variables respond differently to shocks in nominal, floating-rate, and indexed debt. The primary test compares nominal and indexed debt, while the full three-way ranking provides a sharper implication of the valuation mechanism.

The project aims to show that debt composition is not merely a debt-management detail. It provides an empirical window into the mechanism through which fiscal policy affects prices, risk premia, and macroeconomic outcomes.

\appendix

\section{Appendix A: Microfoundation of the Valuation Equation}

This appendix provides the formal derivation of the valuation equation used in the main text.

Consider a representative household with preferences
\begin{equation}
E_0\sum_{t=0}^{\infty}\beta^t U(C_t,M_t/P_t),
\end{equation}
where \(C_t\) is consumption and \(M_t/P_t\) denotes real money balances. The government issues nominal fixed-rate debt \(B_t^N\), floating-rate debt \(B_t^F\), and inflation-indexed debt \(B_t^R\). The household's nominal budget constraint is
\begin{align}
P_t C_t+B_t^N+B_t^F+P_t B_t^R+M_t
=&
P_tY_t+(1+i_{t-1}^N)B_{t-1}^N+(1+i_{t-1})B_{t-1}^F
\nonumber\\
&+P_t(1+i_{t-1}^R)B_{t-1}^R+M_{t-1}-P_tT_t.
\end{align}

Let \(M_{t,t+h}\) denote the real stochastic discount factor. Iterating the government's flow budget constraint forward and imposing the transversality condition yields
\begin{equation}
V_t^N+V_t^F+V_t^R+\frac{M_t}{P_t}
=
E_t\sum_{h=0}^{\infty}M_{t,t+h}(S_{t+h}+\Sigma_{t+h}),
\end{equation}
where \(\Sigma_t\) denotes seigniorage revenue. The main text omits seigniorage for simplicity, yielding equation \eqref{eq:valuation_equation}. Including seigniorage does not alter the ranking logic as long as seigniorage responses are not systematically instrument-specific after controlling for monetary policy.

\section{Appendix B: Construction of Prisma Fiscal Backing Revisions}

Let \(E_t^P[s_{t+q}]\) denote the Prisma Fiscal consensus forecast of the primary surplus for horizon \(q\). The baseline finite-horizon fiscal-backing proxy is
\begin{equation}
\mathcal{F}_t^P
=
\sum_{q=0}^{Q}
\frac{E_t^P[s_{t+q}]}{(1+\bar r)^q}.
\end{equation}
The baseline revision is
\begin{equation}
\Delta \mathcal{F}_t^P
=
\mathcal{F}_t^P-\mathcal{F}_{t-1}^P.
\end{equation}
Robustness checks vary \(Q\), \(\bar r\), and the interpolation method when survey horizons change.

\section{Appendix C: Robustness Region for Residual Backing}

The ranking region in equations \eqref{eq:region_nf}--\eqref{eq:region_fr} follows directly from \(\phi^j=\kappa^j(1-\rho^j)\). For the \(N\)-\(F\) margin,
\begin{align}
\phi^N>\phi^F
&\iff
\kappa^N(1-\rho^N)>\kappa^F(1-\rho^F) \\
&\iff
\kappa^N\rho^N-\kappa^F\rho^F<\kappa^N-\kappa^F.
\end{align}
For the \(F\)-\(R\) margin,
\begin{align}
\phi^F>\phi^R
&\iff
\kappa^F(1-\rho^F)>\kappa^R(1-\rho^R) \\
&\iff
\kappa^F\rho^F-\kappa^R\rho^R<\kappa^F-\kappa^R.
\end{align}
When \(\kappa^R=0\), the second inequality reduces to \(\rho^F<1\). Hence the most fragile margin in the three-way ranking is the \(N\)-\(F\) comparison, because \(\kappa^F\) is regime-dependent.

\begin{thebibliography}{99}

\bibitem[Arellano(2008)]{Arellano2008}
Arellano, C. (2008).
Default risk and income fluctuations in emerging economies.
\emph{American Economic Review}, 98(3), 690--712.

\bibitem[Bi(2012)]{Bi2012}
Bi, H. (2012).
Sovereign default risk premia, fiscal limits, and macroeconomic dynamics.
\emph{European Economic Review}, 56(3), 389--410.

\bibitem[Bianchi and Ilut(2017)]{BianchiIlut2017}
Bianchi, F., and C. Ilut (2017).
Monetary/fiscal policy mix and agents' beliefs.
\emph{Review of Economic Dynamics}, 26, 113--139.

\bibitem[Cochrane(2023)]{Cochrane2023}
Cochrane, J. H. (2023).
\emph{The Fiscal Theory of the Price Level}.
Princeton University Press.

\bibitem[Conley, Hansen, and Rossi(2012)]{ConleyHansenRossi2012}
Conley, T. G., C. B. Hansen, and P. E. Rossi (2012).
Plausibly exogenous.
\emph{Review of Economics and Statistics}, 94(1), 260--272.

\bibitem[Davig and Leeper(2011)]{DavigLeeper2011}
Davig, T., and E. M. Leeper (2011).
Monetary--fiscal policy interactions and fiscal stimulus.
\emph{European Economic Review}, 55(2), 211--227.

\bibitem[Ghosh et al.(2013)]{GhoshEtAl2013}
Ghosh, A. R., J. I. Kim, E. G. Mendoza, J. D. Ostry, and M. S. Qureshi (2013).
Fiscal fatigue, fiscal space and debt sustainability in advanced economies.
\emph{Economic Journal}, 123(566), F4--F30.

\bibitem[Jordà(2005)]{Jorda2005}
Jordà, Ò. (2005).
Estimation and inference of impulse responses by local projections.
\emph{American Economic Review}, 95(1), 161--182.

\bibitem[Leeper(1991)]{Leeper1991}
Leeper, E. M. (1991).
Equilibria under active and passive monetary and fiscal policies.
\emph{Journal of Monetary Economics}, 27(1), 129--147.

\bibitem[Mendoza and Ostry(2008)]{MendozaOstry2008}
Mendoza, E. G., and J. D. Ostry (2008).
International evidence on fiscal solvency: Is fiscal policy responsible?
\emph{Journal of Monetary Economics}, 55(6), 1081--1093.

\bibitem[Olea and Pflueger(2013)]{OleaPflueger2013}
Olea, J. L. M., and C. Pflueger (2013).
A robust test for weak instruments.
\emph{Journal of Business \& Economic Statistics}, 31(3), 358--369.

\bibitem[Ramey and Zubairy(2018)]{RameyZubairy2018}
Ramey, V. A., and S. Zubairy (2018).
Government spending multipliers in good times and in bad: Evidence from US historical data.
\emph{Journal of Political Economy}, 126(2), 850--901.

\bibitem[Romer and Romer(2010)]{RomerRomer2010}
Romer, C. D., and D. H. Romer (2010).
The macroeconomic effects of tax changes: Estimates based on a new measure of fiscal shocks.
\emph{American Economic Review}, 100(3), 763--801.

\bibitem[Sims(1994)]{Sims1994}
Sims, C. A. (1994).
A simple model for study of the determination of the price level and the interaction of monetary and fiscal policy.
\emph{Economic Theory}, 4(3), 381--399.

\bibitem[Woodford(1995)]{Woodford1995}
Woodford, M. (1995).
Price-level determinacy without control of a monetary aggregate.
\emph{Carnegie-Rochester Conference Series on Public Policy}, 43, 1--46.

\end{thebibliography}

\end{document}
